\documentclass[aspectratio=169,11pt]{beamer}
% Shared CMPSC 480 Beamer preamble.
\usepackage[T1]{fontenc}
\usepackage[utf8]{inputenc}
\usepackage{lmodern}
\usepackage{amsmath,amssymb,mathtools}
\usepackage{booktabs,array,tabularx}
\usepackage{xcolor}
\usepackage{listings}
\usepackage{tikz}
\usetikzlibrary{arrows.meta,positioning,shapes.geometric}

% Ignore only negligible TeX box diagnostics that are below visual tolerance.
% Larger layout defects still appear in the build log and in rendered QA.
\vfuzz=3pt
\hbadness=2000

\definecolor{courseink}{HTML}{111111}
\definecolor{coursegray}{HTML}{EDEDED}
\definecolor{courserule}{HTML}{B8BCC4}
\definecolor{courseblue}{HTML}{3D8DFF}
\definecolor{courselightblue}{HTML}{D0EDFA}
\definecolor{coursemuted}{HTML}{5C6470}
\definecolor{coursegreen}{HTML}{0B7A53}
\definecolor{coursered}{HTML}{B3261E}

\usetheme{default}
\usefonttheme{professionalfonts}
\setbeamertemplate{navigation symbols}{}
\setbeamersize{text margin left=0.65cm,text margin right=0.65cm}

\setbeamercolor{normal text}{fg=courseink,bg=white}
\setbeamercolor{frametitle}{fg=courseink,bg=white}
\setbeamercolor{title}{fg=courseink,bg=white}
\setbeamercolor{subtitle}{fg=coursemuted,bg=white}
\setbeamercolor{structure}{fg=courseblue}
\setbeamercolor{alerted text}{fg=coursered}
\setbeamercolor{block title}{fg=courseink,bg=courselightblue}
\setbeamercolor{block body}{fg=courseink,bg=coursegray}
\setbeamercolor{example text}{fg=coursegreen}

\setbeamerfont{title}{size=\fontsize{25}{29}\selectfont,series=\bfseries}
\setbeamerfont{subtitle}{size=\fontsize{13}{16}\selectfont}
\setbeamerfont{frametitle}{size=\fontsize{19}{22}\selectfont,series=\bfseries}
\setbeamerfont{normal text}{size=\fontsize{11}{14}\selectfont}
\setbeamerfont{footline}{size=\fontsize{7.5}{9}\selectfont}

\setbeamertemplate{frametitle}{%
  \nointerlineskip
  % Use content-driven height so a long assessment title can wrap without
  % being clipped by a fixed-height title bar.
  \begin{beamercolorbox}[wd=\paperwidth,sep=0.17cm,leftskip=0.48cm,rightskip=0.48cm]{frametitle}
    \insertframetitle
  \end{beamercolorbox}
  \vspace{-0.08cm}
  {\color{courserule}\rule{\textwidth}{0.35pt}}
}

\setbeamertemplate{footline}{%
  \leavevmode
  \begin{beamercolorbox}[wd=\paperwidth,ht=2.6ex,dp=1.1ex,leftskip=.65cm,rightskip=.65cm]{author in head/foot}
    \color{coursemuted}\insertshorttitle\hfill\insertframenumber/\inserttotalframenumber
  \end{beamercolorbox}
}

\setbeamertemplate{itemize item}{\color{courseblue}\small$\blacktriangleright$}
\setbeamertemplate{itemize subitem}{\color{courseblue}\scriptsize$\bullet$}
\setbeamertemplate{enumerate item}{\color{courseblue}\insertenumlabel.}

\lstdefinestyle{coursepython}{
  language=Python,
  basicstyle=\ttfamily\fontsize{8.5}{10}\selectfont,
  keywordstyle=\color{courseblue}\bfseries,
  commentstyle=\color{coursemuted}\itshape,
  stringstyle=\color{coursegreen},
  showstringspaces=false,
  columns=fullflexible,
  keepspaces=true,
  frame=single,
  rulecolor=\color{courserule},
  backgroundcolor=\color{coursegray},
  breaklines=true,
  tabsize=4,
  xleftmargin=0.15cm,
  xrightmargin=0.15cm,
  framexleftmargin=0.15cm,
  framexrightmargin=0.15cm
}
\lstset{style=coursepython}

\newcommand{\points}[1]{\hfill{\color{courseblue}\textbf{[#1 points]}}}
\newcommand{\deliverable}[1]{\textbf{\color{courseblue}#1}}
\newcommand{\code}[1]{\texttt{#1}}
\newcommand{\keyidea}[1]{%
  \begin{beamercolorbox}[rounded=false,sep=0.55em,wd=\linewidth]{block body}
    \textbf{Key idea.} #1
  \end{beamercolorbox}
}
\newcommand{\answerlabel}{\textbf{\color{coursegreen}Answer.}\ }
\newcommand{\gradinglabel}{\textbf{\color{courseblue}Grading guidance.}\ }

\newenvironment{questionblock}[1]
  {\begin{block}{#1}}
  {\end{block}}

\newenvironment{solutionblock}[1]
  {\begin{block}{#1}}
  {\end{block}}

\AtBeginSection[]{
  \begin{frame}
    \vfill
    {\usebeamerfont{title}\color{courseink}\insertsectionhead\par}
    \vspace{0.4cm}
    {\color{courseblue}\rule{0.32\paperwidth}{3pt}}
    \vfill
  \end{frame}
}


% CMPSC480_TOTAL_POINTS: 10
% CMPSC480_ITEM_IDS: 1,2,3,4
\title[Week 8 Assignment Questions]{Week 8 Assignment: Linear Regression with Batch Gradient Descent}
\subtitle{Raw NumPy, normalization, convergence, and numerical defenses}
\author{CMPSC 480 - Student Question Deck}
\date{}

\begin{document}


\begin{frame}
  \titlepage
\end{frame}

\begin{frame}{Assessment overview}
\begin{columns}[T,onlytextwidth]
\begin{column}{0.62\textwidth}
\Large Build a vectorized regression pipeline and compare it with a stable closed-form oracle.

\vspace{0.35cm}
\normalsize
Coverage: Week 8: linear regression, MSE, gradient descent, and scaling
\end{column}
\begin{column}{0.33\textwidth}
\begin{block}{Points}10 total\end{block}
\begin{block}{Expected time}6-8 hours\end{block}
\begin{block}{Submission}Repository plus convergence report\end{block}
\end{column}
\end{columns}
\end{frame}

\begin{frame}{Learning objectives}
\begin{itemize}
  \item Fit and reuse training-only normalization statistics.
  \item Implement a vectorized MSE gradient.
  \item Handle an intercept consistently.
  \item Stop on bounded, reproducible convergence criteria.
  \item Compare predictions with a pseudoinverse reference.
\end{itemize}
\end{frame}

\begin{frame}{Requirements checklist}
\small
\begin{itemize}
  \item Use Python 3.11 and NumPy; do not use a library that implements the assessed algorithm.
  \item Keep data, model, optimization, and evaluation responsibilities behind explicit interfaces.
  \item Validate shapes, ranges, finite values, empty inputs, and terminal or degenerate cases.
  \item Use deterministic seeds and record every configuration used for reported evidence.
  \item Include focused tests and a README with exact reproduction commands.
\end{itemize}
\end{frame}

\begin{frame}{Task 1 - Data and normalization (2 points)}
\small
Validate X and y and fit z-score scaling on training features.

\vspace{0.2cm}
\begin{enumerate}
\item[\textbf{a.}] Specify fit and predict shape rules and bias handling. \hfill{\color{courseblue}\textbf{[1 points]}}
\item[\textbf{b.}] Handle a zero-variance feature. \hfill{\color{courseblue}\textbf{[1 points]}}
\end{enumerate}
\end{frame}

\begin{frame}{Task 2 - Vectorized optimizer (4 points)}
\small
Implement full-batch gradient descent with cost history and finite guards.

\vspace{0.2cm}
\begin{enumerate}
\item[\textbf{a.}] Derive and implement the gradient of one-half mean squared error. \hfill{\color{courseblue}\textbf{[2 points]}}
\item[\textbf{b.}] Compute one update for X=[[1,0],[1,1]], y=[1,3], w=[0,0], alpha=0.5. \hfill{\color{courseblue}\textbf{[1 points]}}
\item[\textbf{c.}] Define divergence and stopping checks. \hfill{\color{courseblue}\textbf{[1 points]}}
\end{enumerate}
\end{frame}

\begin{frame}{Task 3 - Oracle comparison (2 points)}
\small
Compare learned predictions with NumPy pseudoinverse least squares.

\vspace{0.2cm}
\begin{enumerate}
\item[\textbf{a.}] State why the pseudoinverse is preferable to explicitly inverting X transpose X. \hfill{\color{courseblue}\textbf{[1 points]}}
\item[\textbf{b.}] Define a comparison for a collinear dataset. \hfill{\color{courseblue}\textbf{[1 points]}}
\end{enumerate}
\end{frame}

\begin{frame}{Task 4 - Experiments and tests (2 points)}
\small
Evaluate learning rate and scaling while covering small and degenerate inputs.

\vspace{0.2cm}
\begin{enumerate}
\item[\textbf{a.}] Compare two learning rates with and without scaling. \hfill{\color{courseblue}\textbf{[1 points]}}
\item[\textbf{b.}] Test one feature, m<n, constant feature, empty input, and nonfinite data. \hfill{\color{courseblue}\textbf{[1 points]}}
\end{enumerate}
\end{frame}

\begin{frame}{Edge cases and defensive programming}
\small
\begin{itemize}
  \item A constant feature.
  \item One training example.
  \item More features than examples.
  \item Learning rate causes nonfinite loss.
  \item Predict is called before fit.
\end{itemize}
\end{frame}

\begin{frame}{Submission instructions}
\small
\begin{itemize}
  \item Submit source, tests, README, configuration, and the requested report.
  \item Provide one command that reproduces the primary result from a clean environment.
  \item Exclude credentials, generated caches, and unlicensed data.
\end{itemize}
\vspace{0.2cm}
\alert{Do not submit credentials, generated caches, or unapproved libraries.}
\end{frame}

\begin{frame}{Grading rubric}
\scriptsize
\begin{tabularx}{\textwidth}{>{\bfseries}p{0.28\textwidth} p{0.08\textwidth} X}
\toprule
Category & Points & Required evidence \\
\midrule
Data and normalization & 2 & Stored means/scales and preprocessing tests. \\
Vectorized optimizer & 4 & Trainer, gradient tests, and convergence trace. \\
Oracle comparison & 2 & Tolerance checks on several synthetic datasets. \\
Experiments and tests & 2 & Controlled table and focused test suite. \\
\bottomrule
\end{tabularx}
\end{frame}

\begin{frame}{Reflection prompt}
In 150-250 words:

Explain one case where the loss curve looked plausible but an oracle or invariant exposed a defect.
\end{frame}

\end{document}
